\chapter{Индексный дефект и выбор ориентации цепного KMS-процесса}
\label{ch:v8-chain-orientation-index-defect-selector-gate}

\section{Почему две KMS-ветви нельзя оставить равноправными}

Предыдущая глава получила два операторно допустимых отношения скоростей
$e^{-2}$ и $e^2$, соответствующие $N$ и $2I-N$. Но Том VII начинался не с
произвольной цепи: знак хирального Ходж-сдвига был выбран до вычисления
вакуума и обеспечивал неустойчивость нулевого поля. Поэтому требуется
проверить совместимость обеих ориентаций с тем же рангоизменяющим
родителем.

Полярная коизометрия имеет тип
\begin{equation}
 U:\mathbb C^{11}\longrightarrow\mathbb C^{10},
 \qquad UU^*=I_{10},
 \qquad \dim\ker U=1,
 \qquad \dim\ker U^*=0.
\end{equation}
Следовательно,
\begin{equation}
 \operatorname{ind}U=1.
 \label{eq:v8-polar-fredholm-index}
\end{equation}

\section{Словарь между степенью и хиральной градуировкой}

На концевом носителе введём унаследованную хиральную градуировку
\begin{equation}
 \Gamma=\operatorname{diag}(-I_{11},+I_{10}).
 \label{eq:v8-endpoint-chiral-grading}
\end{equation}
Тогда цепной оператор предыдущей главы удовлетворяет точному тождеству
\begin{equation}
 N_{\partial}-I_{21}=\Gamma,
 \qquad
 (2I_{21}-N_{\partial})-I_{21}=-\Gamma.
 \label{eq:v8-chain-grading-dictionary}
\end{equation}
Кроме того,
\begin{equation}
 -\Tr\Gamma=1=\operatorname{ind}U.
 \label{eq:v8-index-grading-sign}
\end{equation}
Значит, обращение цепи меняет не только направление теплового отношения,
но и знак уже существующего индексного пакета.

\section{Рангоизменяющий тест двух знаков}

Для прямоугольного поля $A:\mathbb C^{11}\to\mathbb C^{10}$ рассмотрим
тот же Ходж-момент, который использовался при допуске родителя:
\begin{equation}
 S_{\pm}(A)=\frac1{21}
 \left\|[d_A,d_A^*]\mp\Gamma\right\|_{\rm HS}^2.
 \label{eq:v8-two-chiral-sign-actions}
\end{equation}
На радиальном пути $A=tU$ прямой знак даёт
\begin{equation}
 S_+(tU)=\frac{1+20(1-t^2)^2}{21}.
 \label{eq:v8-selected-radial-action}
\end{equation}
Нуль стационарен, его полный вещественный гессиан размерности $220$ имеет
собственное значение $-8/21$, а минимум достигается при $t=1$:
\begin{equation}
 S_+(U)=\frac1{21}.
\end{equation}
Именно этот знак порождает коизометрический вакуум и одну защищённую
ядерную линию.

Обращённый знак даёт вместо этого
\begin{equation}
 S_-(tU)=\frac{1+20(1+t^2)^2}{21}.
 \label{eq:v8-reversed-radial-action}
\end{equation}
Его гессиан в нуле равен $+8/21$ на всех $220$ направлениях, а единственный
радиальный минимум находится при $t=0$. Следовательно, обращённая ветвь
уничтожает сам ранговый переход, ради которого был допущен родитель.

\section{Выбранное безразмерное отношение}

Совместимость с индексом $+1$, хиральным знаком и неустойчивостью нуля
оставляет
\begin{equation}
 N_{\partial}=0I_{11}\oplus2I_{10},
 \qquad
 \frac{\gamma_\uparrow}{\gamma_\downarrow}=e^{-2}.
 \label{eq:v8-index-selected-rate-ratio}
\end{equation}
Соответствующая направленная KMS-полугруппа из предыдущего гейта имеет
одномерную неподвижную алгебру и щель
$0.0212599674\ldots$. Выбор не использует непрерывной подгонки или
феноменологических данных.

Однако результат остаётся безразмерным. Индекс выбирает ориентацию и
отношение прямой и обратной скоростей, но не задаёт общие интенсивности
связывающих, $QLYR$-, $XLdR$- и калибровочных слагаемых и не переводит параметр
полугруппы в физическое время.

\begin{theorem}[Индексный выбор цепной ориентации]
В классе двух цепных KMS-ориентаций только
$N_{\partial}-I_{21}=\Gamma$ совместимо с унаследованным индексом
$\operatorname{ind}U=1$ и рангоизменяющим Ходж-функционалом. Обращённый
знак делает нулевое поле строго устойчивым. Поэтому безразмерное отношение
скоростей однозначно равно $e^{-2}$ в пределах данного родителя.
\end{theorem}

\begin{nogocriterion}[Запрет превращать индексное отношение в абсолютную скорость]
Число $e^{-2}$ связывает только две противоположные скорости одной
боровской пары. Оно не фиксирует их общий множитель, относительные
интенсивности разных семейств, единицу энергии или физическое время.
\end{nogocriterion}

\begin{protocol}
\begin{enumerate}
 \item индекс полярной коизометрии: \textbf{$+1$};
 \item центрированная цепная степень: \textbf{совпадает с $\Gamma$};
 \item прямой нулевой гессиан: \textbf{$(-8/21)I_{220}$};
 \item прямой минимум: \textbf{коизометрия, энергия $1/21$};
 \item обращённый нулевой гессиан: \textbf{$(+8/21)I_{220}$};
 \item обращённый минимум: \textbf{нулевое поле};
 \item ориентационная двузначность: \textbf{закрыта};
 \item выбранное KMS-отношение: \textbf{$e^{-2}$};
 \item непрерывная подгонка: \textbf{не использована};
 \item базовая метрика интенсивностей: \textbf{не выведена};
 \item абсолютная скорость: \textbf{не выведена};
 \item следующий тест: \textbf{единая цепная метрика форм Дирихле}.
\end{enumerate}
\end{protocol}

Машинный сертификат сохранён в
\path{s2t_v8_chain_orientation_index_defect_selector_gate_results.json}.

\begin{thebibliography}{9}
\bibitem{v8-index-gesztesy}
F.~Gesztesy, Y.~Latushkin, K.~A.~Makarov, F.~Sukochev, Y.~Tomilov,
\emph{The index formula and the spectral shift function for relatively
trace class perturbations}, Advances in Mathematics \textbf{227} (2011),
319--420; arXiv:1004.1582.
\bibitem{v8-index-project}
Том VII, гейты хиральной неустойчивости Ходжа, полярной коизометрии и
цепного относительного факторпространства.
\end{thebibliography}